
\documentclass[12pt,a4paper]{article}
\usepackage[utf8]{inputenc}
\usepackage[english,russian]{babel}
\usepackage{indentfirst}
\usepackage{misccorr}
\usepackage{graphicx}
\usepackage{amsmath}

\graphicspath{{pictures/}}
\DeclareGraphicsExtensions{.pdf,.png,.jpg, .mps, }


\begin{document}






\begin{titlepage}

 \begin{center}
 Московский Государственный Университет имени М.В.Ломоносова \\
Механико-математический факультет\\
Кафедра теории вероятностей
  \end{center}

 \vspace{3cm}
 
 \begin{center}
   
  {  Курсовая работа за 3 курс:\\
  Принципы назначения страховых премий и меры риска}
   
    \vspace{5cm}
\end{center}     
   
     
   \hspace{170pt}  {Выполнила: Александра Токаева,  309 \\}
       
 \vspace{0.1cm}
  \hspace{170pt} 	  Научный руководитель:  проф. Г.И.Фалин\\

\vspace{4cm}

  \begin{center}
  {Москва\\
  2020}
  \end{center}  
  
\newpage
\tableofcontents
 
 
\end{titlepage}

\section{ От автора}


Задачи назначения страховых премий и определения оптимальной цены для различных финансовых инструментов играют важнейшую роль в страховой 
математике, поскольку без правильно назначенной цены на продукт не представляется возможным дальнейший анализ  и оптимизация деятельности страховой или инвестиционной компании. Несмотря на относительную простоту формулировки  - определить оптимальную цену, решение данной задачи требует развития специального математического аппарата  мер риска, сопоставляющих риску числовой (часто денежный ) эквивалент, который позволяет сравнивать риски и выбирать наиболее предпочтительный вариант. Цель данной работы - подробно  и основательно изучить этот раздел современной теории страхования,   дать описание  основных подходов к решению поставленной задачи и важнейших инструментов, используемых в работе, а также сделать определенные выводы. Наши рассуждения в основном опираются на статьи G.I.Falin, On the optimal pricing of a portfolio of heterogeneous insurance risks, 2008,MSU и Landsman, Sherris, Risk measures and Insurance premium principles, Insurance: Mathematics and Economics. Отметим, что  некоторые рассуждения  и логические переходы в вышеназванных монографиях содержат пропуски или вовсе опущены. Мы полностью восстановим все пропущенные рассуждения и добавим важные, на наш взгляд, детали. К таким мы относим, например, альтернативное доказательство задачи 2 (посвященной поиску оптимальных премий, которые  минимизируют взвешеную сумму квадратов разностей между индивидуальными премиями и индивидуальными выплатами и  использующее  не неравенство Коши-Буняковского-Шварца, а минимизационные  свойства перпендикуляра, опущенного из точки на гиперплоскость), доказательство теоремы о существовании функции полезности, согласованной с отношением предпочтения, а также комментарии к использованию метода множителей Лагранжа. Однако мы не претендуем на авторство конкретных утверждений и результатов, а также используемых понятий из теории вероятностей и страхования, поэтому вся работа, проделанная лично нами, будет отдельно отмечена.  Кроме того, для более детального восприятия, в работу добавлены графики, самостоятельно выполненные  нами в средах программирования Python и Metapost, а также приведен код, необходимый для их создания. \\
В первой части  мы применяем самый естественный подход, когда в качестве числовой характеристики риска используется его математическое ожидание.
Мы  рассматриваем наборы  индивидуальных премий как векторы в многомерном евклидовом пространстве и  используем  геометрические принципы и факты из многомерной геометрии, чтобы найти оптимальные значения этих премий и минимизировать вероятность  разорения. \\
Вторая часть дает описание других  подходов к назначению числового эквивалента риску: через отношение предпочтения и  функции полезности,  замену вероятностной меры,  функции искажения, а также комбинацию этих подходов.  Наконец, обсуждаются сильные и слабые стороны этих подходов.


\part{ Минимизационные задачи}






{\center Мы применим простые геометрические принципы, чтобы показать, что хорошо известные подходы к назначению премий 
в страховом контракте минимизируют взвешенные квадраты разностей между индивидуальными премиями и индивидуальными выплатами, а также между суммарными премиями для классов однородных рисков и суммарными выплатами, поступающими из них. }\\


{   \section{ Постановка задачи }} 
 
Рассмотрим портфель из $n$ неоднородных независимых страховых рисков. Пусть $X_i$ обозначает размер выплат по$ i-$му риску за рассматриваемый период, S - суммарные потери, связанные с портфелем. При некоторых естественных предположениях (что портфель достаточно большой, не очень неоднородный и распределение размера выплат не очень ассиметричное) распределение случайной величины $\frac{S-ES}{\sqrt{VarS}}$ может быть приближено стандартным гауссовским распределением.\\
Предположим, что страховщик взимает премию $\pi_i$ по $i-$му риску и таким образом собирает суммарную премию
$\pi=\sum \limits_{i=1}^{n}\pi_i$. Из гаусовости  распределения  величины $\frac{S-ES}{\sqrt{VarS}}$ получаем, что для гарантии достаточно маленькой вероятности разорения $R=P(S>\pi)$ (например, R=5\%) страховщик должен собрать суммарную премию в размере 
$$ES + \sqrt{VarS}*z_{(1-R)} \eqno(1) $$, где $z_{\alpha} $
 - квантиль гаусовского распределения уровня альфа.
 
 \begin{figure}[h]
\center{\includegraphics[scale=0.9]{pictures1.mps}}
\caption{ Квантиль уровня $\alpha $}
\label{fig:image}
\end{figure}




Ну действительно: $P(S>\pi)=R \Leftrightarrow P(\frac{S-ES}{\sqrt{VarS}} > \frac{\pi-ES}{\sqrt{VarS}})=R \Leftrightarrow \frac{S-ES}{\sqrt{VarS}}=z_{(1-R)} \Leftrightarrow \pi=ES + \sqrt{VarS}*z_{(1-R)}.$\\
Последнее равенство ничего не говорит о величине индивидуальных премий. Чтобы найти их, мы используем дополнительные принципы. \\
Вслед за Заксом, Фростигом и Левиксоном мы рассмотрим два подхода к задаче разбиения величины $\pi$ на n индивидуальных премий $\pi_1 \ldots \pi_n:$\\

1) Для заданной вероятности разорения $R=P(S>\pi)$ минимизировать взвешенный квадрат разности $\sum \limits_{i=1}^{n} \frac {1}{s_i} E(X_i-\pi_i)^2$ между индивидуальными рисками $X_i$ и индивидуальными премиями $\pi_i$ (где $s_i$ -это некоторые известные положительные числа, то есть веса)\\

2) Для заданной $D=\sum \limits_{i=1}^{n} \frac {1}{s_i} E(X_i-\pi_i)^2$  минимизировать вероятность разорения $P(S>\pi)$\\



Сейчас мы с помощью простых геометрических рассуждений покажем, что обе задачи минимизации имеют одно и то же решение. Кроме того, мы покажем, что оптимальные премиии $\pi_i$ минимизируют взвешенную сумму  квадратов разностей между индивидуальными премиями и индивидуальными выплатами, а также между суммарными премиями для классов однородных рисков и суммарными выплатами, поступающими из них. \\
Наше достижение опирается на недавнюю статью, написанную Заксом, Фростигом и Левиксоном, которые исследовали похожие задачи оптимальных цен на  неоднородный портфель( который может быть разделен на классы однородных рисков) с помощью алгебраических методов, основанных на теоремах о положительно определенных матрицах.\\


{\section { Общие результаты о случайных величинах}}
{\subsection{  Проблема оптимизации}}

Пусть $\xi_1, \ldots ,\xi_N$ - случайные величины с конечными матожиданиями $a_1 \cdots a_N$ и дисперсиями $Var\xi_1 \ldots Var\xi_N$. Мы предполагаем, что матожидание и дисперсии известны.\\
Нам бы хотелось заменить случайные величины $\xi_1, \ldots , \xi_N$ на неслучайные числа $A_1, \ldots, A_n$ таким образом, чтобы взвешенная сумма 
$$D \equiv \sum \limits_{i=1}^{n} \omega_i E(\xi_i-A_i)^2 \eqno(2)$$
 была бы минимальна.
Здесь $\omega_i, \ldots, \omega_N$ - это известные числа(веса). \\
Используя элементарные свойства случайных величин, мы можем переписать $D$ следующим образом:
$$D= \sum \limits_{i=1}^{n} \omega_i E(\xi_i-A_i)^2 = \sum \limits_{i=1}^{n} \omega_i( Var(\xi_i-A_i) + (E(\xi_i-A_i))^2)$$ = $$ \sum \limits_{i=1}^{n} \omega_i( Var(\xi_i-A_i) + (a_i-A_i)^2)= \sum \limits_{i=1}^{n} \omega_i Var\xi_i + \sum \limits_{i=1}^{n} \omega_i( a_i-A_i)^2 \eqno(3)$$\\

Поскольку $\omega_i$ и $Var\xi_i$ фиксированы, то изначальная задача минимизации превращается в задачу нахождения минимального значения функции
$$f(A_1, \ldots, A_N) = \sum \limits_{i=1}^{n} \omega_i( a_i-A_i)^2 \eqno(4)$$
Очевидно, оптимальным значением являются $$A_1^*=a_1, \ldots, A_N^*=a_N$$
и минимальное значение этой функции равно нулю. Соответственно, минимальное значение величины $D$ равно $\sum \limits_{i=1}^{n} \omega_i Var\xi_i $\\
Более интересной задача становится, если мы  накладываем  дополнительные ограничения на переменные $A_1, \ldots, A_N$. Принимая во внимания последующее приложение этой задачи к страхованию, мы рассматриваем следующую задачу:\\

{\bf Задача 1 }\\ Найти минимальное значение $D$ при условии, что $$A_1 + \ldots+  A_N=C \eqno(5)$$, где C-известная константа.

Благодаря (3), достаточно найти минимальное значение функции(4) на множестве(5).\\

Для решения этой задачи введем новые переменные $x_i=\sqrt \omega_i (A_i-a_i)$, то есть 
$A_i=a_i + \frac{1} {\sqrt {\omega_i}} x_i$. Тогда задача 1 превращается в:\\

{\bf Задача 2 }\\Найти минимальное значение функции $$g(x_1, \ldots, x_N)= \sum \limits_{i=1}^{n} x_i^2 \eqno(6)$$ при условии, что 
$$\sum \limits_{i=1}^{n} \frac{1}{\sqrt \omega_i} x_i = C - \sum \limits_{i=1}^{n} a_i \eqno(7) $$.\\

Последовательности $X=(x_1, \ldots, x_N)$ и $Y=(\frac{1}{\sqrt {\omega_1} }, \ldots, \frac{1}{\sqrt {\omega_N}}$ можно понимать как N-мерные евклидов векторы в пространстве $R^N$. Соответственно, левая часть равенства(7) есть скалярное произведение X и Y, а функция $g(x_1, \ldots, x_N)$ есть $||X||^2$, где $$||X||= \sqrt{x_1^2+
\ldots+x_N^2}$$ - это длина вектора X.\\

Последующие рассуждения основаны на неравенстве Коши-Буняковского-Шварца, согласно которому для любых двух векторов $X,Y \in R^N$ верно 
$$|X \cdot Y| \leq ||X|| \cdot ||Y||$$ , причем равенство достигается тогда и только тогда, когда X и Y линейно зависимы(в частности, если вектор Y ненулевой, линейная зависимость означает, что X пропорционален Y: $ X=t \cdot Y$ для некоторого $t \in R$)\\

Применяя это неравенство, получаем:
$$g(x_1, \ldots, x_N)=||X||^2 \geq \frac{|X \cdot Y|^2}{||Y||^2} = \frac{(C-\sum\limits_{i=1}^{n} a_i)^2}
{\sum\limits_{i=1}^{n} \frac{1}{\omega_i} }  \eqno(7.5)$$

Поэтому для $(x_1, \ldots, x_N)$, удовлетворяющих (7), имеем:
$$ min g(x_1, \ldots, x_N) \geq \frac{(C-\sum\limits_{i=1}^{n} a_i)^2}{\sum\limits_{i=1}^{n} \frac{1}{\omega_i}} 
\eqno(8) $$

Поскольку вектор $Y=(\frac{1}{\sqrt{\omega_1}}, \ldots, \frac{1}{\sqrt{\omega_N}})$ ненулевой, то равенство в (8) достигается тогда и только тогда когда существует такое t что 
$$x_i = \frac{1}{\sqrt{\omega_i}} \cdot t , i=1, \ldots, N  \eqno(9)$$

Подставляя, что $X=t \cdot Y$ в (7.5), получаем, что 
$$ t^* = \frac{ C-\sum\limits_{i=1}^{n} a_i} {\sum\limits_{i=1}^{N} \frac{1}{\omega_i}}$$, 

$$ x_i^*=\frac{1}{\sqrt{\omega_i}} \cdot t^* $$

$$A_i^*=a_i + \frac{1}{\sqrt \omega_i x_i^*} = a_i + \frac{1}{\omega_i}t^*
 = a_i + \frac{1}{\omega_i} \frac{C-\sum\limits_{j=1}^{N} a_j}{\sum\limits_{j=1}^{N} \frac{1}{\omega_j} } \eqno(10)$$

$$D_{min}= \sum \limits_{i=1}^{N} \omega_i Var\xi_i  +  \frac{(C-\sum\limits_{j=1}^{N} a_j)^2}
{\sum\limits_{j=1}^{N} \frac{1}{\omega_j}} \eqno(11)$$


{\subsection{  Альтернативное решение задачи 2}}



Дадим альтернативное решение задачи $2$ :\\



{\bf Задача 2 }\\Найти минимальное значение функции $$g(x_1, \ldots, x_N)= \sum \limits_{i=1}^{n} x_i^2 $$ при условии, что 
$$\sum \limits_{i=1}^{n} \frac{1}{\sqrt \omega_i} x_i = C - \sum \limits_{i=1}^{n} a_i $$.\\

Заметим, что наборы чисел $X=(x_1, \ldots, x_N)$ и $Y=(\frac{1}{\sqrt {\omega_1} }, \ldots, \frac{1}{\sqrt {\omega_N}} )$ можно понимать как  $N$-мерные евклидовы векторы в пространстве $R^N$. Соответственно, левая часть  последнего равенства есть скалярное произведение X и Y, а функция $g(x_1, \ldots, x_N)$ есть $||X||^2$, где $$||X||= \sqrt{x_1^2+
\ldots+x_N^2}$$ - это длина вектора X.\\

То есть наша задача заключается в том, чтобы  минимизировать квадрат длины вектора $X, $ удовлетворяющего условию 
$\sum \limits_{i=1}^{n} \frac{1}{\sqrt \omega_i} x_i = C - \sum \limits_{i=1}^{n} a_i $. Но заметим, что данное условие означает, что вектор $X$ принадлежит гиперплоскости с нормалью $Y=(\frac{1}{\sqrt {\omega_1} }, \ldots, \frac{1}{\sqrt {\omega_N}} )$ .\\

Последнее утверждение требует некоторых пояснений. Как известно, гиперплоскость- это линейная поверхность коразмерности один, то есть линейная оболочка $ n-1 $вектора.  Из линейной алгебры известно, что линейные пространства можно задавать системами линейных уравнений, причем если система имеет ранг $ k$, то задаваемое ей пространство будет иметь размерность $n-k$.  (это доказано на стр. 4  книги Шурыгина В.В, "Аналитическая геометрия, часть 3" ). В нашем случае размерность пространства равна $n-1$, поэтому для задания  гиперплоскости в $n$-мерном пространстве  требуется всего одно уравнение. Запишем его в виде $a_1x_1 + \dots + a_nx_n=c$. Согласно общей теории, это уравнение задает плоскость размерности n-1. Но с другой стороны, это уравнение можно переписать в виде  $ <a,x> =c$, то есть вектор a перпендикулярен всем векторам x из этой гиперплоскости.  То есть a - вектор нормали к данной гиперплоскости.\\

Данное утверждение, сформулированное как  " в ортонормированной  системе координат главный вектор плоскости является и нормальным ее вектором" доказано на стр. 99  учебника A.E.Умнова "Аналитическая геометрия и линейная алгебра", Москва, МФТИ, 2011.\\

 \begin{figure}[h!]
\center{\includegraphics[scale=0.9]{mproof.pdf}}
\caption{ Иллюстрация}
\label{fig:image}
\end{figure}




Из курса линейной алгебры известно, что минимизирует расстояние от точки до гиперплоскости - перпендикуляр, опущенный из этой точки на гиперплоскость (это непосредственно следует из многомерной  теоремы Пифагора).\\
 Но перпендикуляр - это вектор, параллельный нормали к поверхности, а  выше мы уже пояснили, что нормаль к  нашей гиперплоскости - это вектор $Y$.\\
 
 
$\Rightarrow X=tY,$ то есть $(x_1, \dots, x_n) = t(\frac{1}{\sqrt {\omega_1} }, \ldots, \frac{1}{\sqrt {\omega_N}} )$\\

Подставим выражение для X в условие  $(X,Y) = C- \sum \limits_{i=1}^{n} a_i $\\

Получим $t||Y||^2 = C -  \sum \limits_{i=1}^{n} a_i $\\



$ \Rightarrow  t^* =   \frac{ C-\sum\limits_{i=1}^{n} a_i} { ||Y||^2} = \frac{ C-\sum\limits_{i=1}^{n} a_i} {\sum\limits_{i=1}^{N} \frac{1}{\omega_i}}$, 


$\Rightarrow min  ||X||^2 =  t^2 ||Y||^2 = t(t||Y||^2) = t(C - \sum \limits_{i=1}^{n} a_i ) =  \frac{ (C-\sum\limits_{i=1}^{n} a_i )^2} {\sum\limits_{i=1}^{N} \frac{1}{\omega_i}}$


$\Rightarrow D_{min}= \sum \limits_{i=1}^{N} \omega_i Var\xi_i  +  \frac{(C-\sum\limits_{j=1}^{N} a_j)^2}
{\sum\limits_{j=1}^{N} \frac{1}{\omega_j}} $


\begin{figure}[h]
\caption{ }\label{pic-jan21-2008-01} \unitlength=1.4mm
\begin{picture}(100,70)
\put(10,0){\vector(0,1){60}} \put(0,10){\vector(1,0){90}}
\put(4,6){\vector(3,2){60}}
\put(92,10){$x_1$}\put(10,62){$x_3$}\put(67,45){$x_2$}\put(10,10){\circle*{1}}
\put(11,6){$O$} \qbezier(10,50)(10,50)(80,10)
\qbezier(10,50)(10,50)(55,40)\qbezier(55,40)(55,40)(80,10)
\put(10,10){\line(2,3){21}}\put(31,41.5){\circle*{1}}\put(32,38){$(x_1^*,x_2^*,x_3^*)$}
\put(10,10){\line(3,1){50}}\put(60,26.6){\circle*{1}}\put(52,28){$(x_1,x_2,x_3)$}
\put(34,15){distance=$g(x_1,x_2,x_3)$}
\put(70,25){hyperplane $\sum_{i=1}^{N}\frac{1}{\sqrt{w_i}}x_i=C'$}
\end{picture}
\end{figure}







{\subsection { Двойственная проблема оптимизации}}
Такой же подход может быть применен для изучения двойственной задачи оптимизации:\\
{\bf Задача 3} Найти максимум суммы $A_1 + \ldots + A_N$ если  задано 
$$D= \sum \limits_{i=1}^{n} \omega_i E( \xi_i-A_i)^2 \eqno(12)$$


Заметим, что из (3) следует, что константа D должна быть больше или равна чем $\sum \limits_{i=1}^{n} \omega_i Var\xi_i$\\

Как и раньше, перепишем $$ \sum \limits_{i=1}^{n} \omega_i E( \xi_i-A_i)^2  = \sum \limits_{i=1}^{n} \omega_i Var\xi_i + \sum \limits_{i=1}^{n} \omega_i( a_i-A_i)^2 $$

Ограничение (12) превращается в равенство $$ D^{'} = \sum \limits_{i=1}^{n} \omega_i( a_i-A_i)^2 $$, где  
$$D^{'} = D - \sum \limits_{i=1}^{n} \omega_i Var\xi_i \geq 0 $$

Вводя $x_i=\sqrt \omega_i (A_i-a_i)$, мы сводим задачу 3 к следующему виду:\\


{\bf Задача 3} Найти максимум суммы $\sum \limits_{i=1}^{n} \frac{1}{\sqrt \omega_i} x_i$ если  задана
сумма $$D^{'} = \sum \limits_{i=1}^{N} x_i^2 \eqno(13)$$

Аналогично задаче 2, применяем неравенство Коши-Буняковского-Шварца:\\
$$\sum \limits_{i=1}^{N} \frac{1}{\sqrt \omega_i} x_i = X \cdot Y \leq ||X|| \cdot ||Y|| = \sqrt{D^{'}} \sqrt{\sum\limits_{i=1}^{N} \frac{1}{\omega_i^2}} \eqno(14)$$

Равенство в (14) достигается тогда и только тогда, когда существует t такое что 
$$x_i = \frac{1}{\sqrt{\omega_i}} \cdot t , i=1, \ldots, N  \eqno(15)$$

Подставляя выражение $X=t \cdot Y$ в (14), получаем единственное решение 
$$ t^* = \sqrt{ \frac{D^{'}}{\sum\limits_{i=1}^{N} \frac{1}{\omega_i^2}}}$$
$$x_i^* = \frac{1}{\sqrt{\omega_i}} \cdot t^* = \frac{1}{\sqrt{\omega_i}} \sqrt{ \frac{D^{'}}{\sum\limits_{i=1}^{N} \frac{1}{\omega_i^2}}}$$


Соответственно, такие $A_i$ задают решение оптимизационной задачи 3:
$$A_i^* = a_i + \frac{1}{\sqrt{\omega_i}} \cdot t^* = \frac{1}{\sqrt{\omega_i}} \sqrt{ \frac{D^{'}}{\sum\limits_{i=1}^{N} \frac{1}{\omega_i^2}}} \eqno(16)$$



{\section{ Приложение к модели индивидуального риска}}

В этом разделе мы применим полученные выше результаты к задаче оптимального назначения премий для неоднородного портфеля, рассмотренной Заксом, Фростигом и Левиксоном.\\

Рассмотрим модель индивидуального риска:\\
$$S= X_1 + \ldots + X_n,$$ где n - это общее число рисков в портфеле, случайная величина $X_i$ обозначает потери по i-my риску за рассматриваемый период, а S - это общие потери по портфелю.\\

Мы предполагаем, что случайные величины $X_1, \ldots, X_n$ независимы и имеют конечные матожидание $\mu_1, \ldots, \mu_n$ и дисперсии $\sigma_1^2, \ldots, \sigma_n^2$ соответственно. Тогда случайная величина S имеет конечные матожидание $\mu = \mu_1+ \ldots + \mu_n$ и дисперсию $\sigma^2 = \sigma_1^2+ \ldots + \sigma_n^2.$  Мы также предполагаем, что для достаточно больших n функция распределения центрированной и нормированной величины полных потерь $\frac {S-\mu}{\sigma}$ может быть приближена функцией распределения стандартной гауссовской величины 
$\Phi(x) =  \int \limits_{\infty}^{\infty}exp^{-\frac{t^2}{2}} dt$, то есть:

$$P(\frac{S-\mu}{\sigma} < x) \approx \Phi(x)$$


Предположим, что страховщик взимает премию $\pi_i$ по i-му риску, то есть всего собирает сумму $\pi=\sum \limits_{i=1}^{n}\pi_i.$ Тоогда вероятность разорения ( это вероятность того, что S будет больше собранной суммы $\pi$ дается формулой $ R=P(S>\pi).$\\
Используя гаусовость  $\frac{S-\mu}{\sigma}$, получаем, что:
$$R=P(\frac{S-\mu}{\sigma} > \frac{\pi-\mu}{\sigma} )\approx 1 - \Phi(\frac{\pi-\mu}{\sigma}). \eqno(17)$$

Предположим, что страховщик готов принять достаточно маленький риск разорения R (например, $R=1\%$). Тогда равенство (17) дает следующую(приближенную) формулу для суммарной премии:

$$\pi = \mu  + \sigma \cdot z_{(1-R)} \eqno(1) \eqno(18) $$, где $z_{\alpha} $
 - квантиль гаусовского распределения уровня альфа, то есть $\Phi(z_{\alpha}) = \alpha$.\\

Равенство (18) ничего не говорит про величины индивидуальных премий $\pi_i.$ Чтобы найти их, нам придется применить дополнительные принципы.

{\subsection { Минимизация разности между индивидуальными рисками и индивидуальными  премиями при заданной вероятности разорения }}

Рассмотрим взвешенную сумму $$ D= \sum \limits_{i=1}^{N} \frac{1}{s_i} E( X_i-\pi_i)^2$$
между индивидуальными рисками $X_1, \ldots, X_n$ и индивидуальными премиями $\pi_1, \ldots, \pi_n$ (где $s_1,   
\ldots, s_n$ - это некие известные положительные числа(веса)) и найдем минимум D:

$$ D \equiv D(\pi_1, \ldots, \pi_n) \rightarrow min. \eqno(19)$$

Применяя формулу(10) для $N=n, \xi_i=X_i, a_i = \mu_i, A_i=\pi_i, \omega_i=\frac{1}{s_i}, C= \mu + \sigma \cdot z_{(1-R)}$ мы можем утверждать, что минимизационная задача (19) с ограничением (18) имеет единственное решение 

$$ \pi_i^* = \mu_i +  \frac{s_i}{\sum\limits_{j=1}^{n} s_j } \cdot \sigma \cdot z_{(1-R)}. \eqno(20)$$


Теперь рассмотрим постановку задачи, рассмотренной Заксом, Фростигом и Левиксоном. Пусть портфель можно разбить на к классов однородных рисков с одинаковыми статистическими характеристиками потерь (обычно риски из одного класса принадлежат одному и тому же сектору бизнеса). Пусть $i-$й класс состоит из $n_i$ рисков с одинаковым средним $\mu_i$
и одинаковыми дисперсиями $\sigma_i^2.$ Тогда общее количество потерь $S_i$  в i-м классе  имеет среднее значение $ES_i= n_i \mu_i$ и дисперсию $ Var S_i = n_i \sigma_i^2.$ Общее число потерь от всего портфеля есть 
$S= S_1 + \ldots + S_k$, причем $\mu \equiv ES= \sum\limits_{i=1}^{k} n_i \mu_i $ , 
$\sigma^2 \equiv VarS= \sum\limits_{i=1}^{k} n_i \sigma_i^2 $


Благодаря однородности рисков внутри отдельного класса i, страховщик должен взимать со всех рисков в этом классе одну и ту же премию $\pi_i.$ Тогда общая премия за все риски в портфеле равна $\pi = \sum\limits_{i=1}^{k} n_i \pi_i.$\\

Рассмотрим $$  D= \sum \limits_{i=1}^{k} \frac{1}{r_i} E( S_i-n_i \pi_i)^2$$ между суммарными потерями по разным секторам бизнеса $S_1, \ldots, S_k$ и суммарными премиями $n_1 \pi_1, \ldots, n_k \pi_k$ от этих классов (где $r_1, \ldots, r_k$ - это некоторые известные положительные числа) и минимизируем D:

$$D = D \equiv D(\pi_1, \ldots, \pi_k) \rightarrow min. \eqno(21)$$


Для того, чтобы получить предписанную вероятность разорения, нужно, чтобы (18)  выполнялось.\\

Применяя формулу(10) для $N=k, \xi_i=S_i, a_i = n_i \mu_i, A_i=n_i \pi_i, \omega_i=\frac{1}{r_i}, C= \mu + \sigma \cdot z_{(1-R)}$ мы можем утверждать, что минимизационная задача (19) с ограничением (18) имеет единственное решение 

$$ n_i \pi_i^* = n_i \mu_i +  \frac{r_i}{\sum\limits_{j=1}^{k} r_j } \cdot \sigma \cdot z_{(1-R)}
\Leftrightarrow  \pi_i^* = \mu_i +  \frac{r_i}{n_i \sum\limits_{j=1}^{k} r_j } \cdot \sigma \cdot z_{(1-R)}. \eqno(22)$$

Теперь вернемся к минимизационной задаче (19) c ограничением (18) и положим для всех рисков из $i-$го класса одинаковое значение параметра $s$ равным $\frac {r_i}{n_i}.$ Тогда  из (20) видно, что оптимальное решение для минимизационной задачи (19) совпадает с оптимальным решением минимизационной задачи (21). Таким образом, одни и те же значения премий минимизируют взвешенную сумму квадратов разностей как между индивидуальными премиями и индивидуальными потерями, так и между суммарными премиями для классов однородных рисков и суммарными потерями для этих блоков.


{\subsection {  Минимизация вероятности разорения при заданной разности между индивидуальными рисками и индивидуальными  премиями }}


{\bf Задача 5} Для модели индивидуального риска $$S= X_1 + \ldots + X_n$$ минимизировать вероятность разорения $R=P(S > \pi)$ при заданной величине $$D = \sum\limits_{i=1}^{n} \frac{1}{s_i} E(X_i - \pi_i)^2$$

Поскольку $P(S > \pi)$ уменьшается при увеличивающемся $\pi$, то задача состоит в нахождении максимального значения суммарной премии $\pi= \pi_1+ \ldots + \pi_n.$\\

Применяя формулу(16) для $N=n,  \xi_i = X_i, a_i= \mu_i, A_i= \pi_i , \omega_i= \frac{1}{s_i},$  мы можем утверждать, что минимизационная задача 5  имеет единственное решение 

$$\pi_i^* = \mu_i + s_i  \sqrt{ \frac{D - \sum\limits_{i=1}^{n} \frac{1}{s_i} \sigma_i^2 }{\sum\limits_{i=1}^{n} s_i}} \eqno(23)$$


Теперь опять предположим, что портфель может быть разделен на k классов однородных рисков с одинаковыми статистическими свойствами потерь. Пусть i-й класс состоит из $n_i$ рисков с одинаковым средним $\mu_i$
и одинаковыми дисперсиями $\sigma_i^2.$ Тогда общее количество потерь $S_i$  в i-м классе  имеет среднее значение $ES_i= n_i \mu_i$ и дисперсию $ Var S_i = n_i \sigma_i^2.$ Общее число потерь от всего портфеля есть 
$S= S_1 + \ldots + S_k$, причем $\mu \equiv ES= \sum\limits_{i=1}^{k} n_i \mu_i $ , 
$\sigma^2 \equiv VarS= \sum\limits_{i=1}^{k} n_i \sigma_i^2 $\\
 
Благодаря однородности рисков внутри отдельного класса i, страховщик должен взимать со всех рисков в этом классе одну и ту же премию $\pi_i.$ Тогда общая премия за все риски в портфеле равна $\pi = \sum\limits_{i=1}^{k} n_i \pi_i.$\\

Рассмотрим оптимизационную задачу:

{\bf Задача 6} Минимизировать вероятность разорения $R=P(S>\pi)$ при заданной величине 
$$D = \sum\limits_{i=1}^{k} \frac{1}{r_i} E(S_i - n_i \pi_i)^2$$

Поскольку $P(S > \pi)$ уменьшается при увеличивающемся $\pi$, то задача состоит в нахождении максимального значения суммарной премии $\pi= n_1 \pi_1+ \ldots + n_k \pi_k.$\\

Применяя формулу(16) для $N=k, \xi_i=S_i,  a_i= n_i \mu_i, A_i= n_i \pi_i , \omega_i= \frac{1}{r_i},$  мы можем утверждать, что минимизационная задача 6  имеет единственное решение 

$$\pi_i^* = \mu_i + \frac{r_i}{n_i}  \sqrt{ \frac{D - \sum\limits_{i=1}^{k} \frac{1}{r_i} n_i \sigma_i^2 }{\sum\limits_{i=1}^{k} r_i}} \eqno(24)$$











Таким образом, мы поняли, как назначать страховые премии исходя из минимизации риска разорения, а теперь разработам теорию меры риска.\\

\part {Меры риска}


В актуарной науке недавно появились меры риска, основанные на искаженных вероятностях, которые используются 
для формирования страховых тарифов. Мера риска должна удовлетворять свойствам  неприятия риска  и  предпочтения разнообразия ,  быть аддитивной, а также работать одинаково  как для  страховых рисков, так и  инвестиционных.  Мера риска, основанная на специальной замене меры, удовлетворяет этим свойствам. Мера риска, основанная на искаженных вероятностях, не всегда аддитивна, а также непоследовательна в своем обращении со страхованием и инвестиционными рисками.  Мы предлагаем еще одну меру риска, которая обладает свойствами неприятия риска и предпочтения разнообразия, является аддитивной и  последовательной в своем отношении к страхованию и инвестиционным рискам. \\

\section{ Постановка задачи}
В актуарной науке модели используются как для количественной оценки рисков, так и для ценообразования. 
В задаче количественной оценки рисков требуется мера риска для преобразования случайной будущей прибыли или убытка  в  детерминированный
эквивалент , который может потом  использоваться для упорядочивания рисков и для целей принятия решений. Для количественной оценки рисков требуется указать распределение вероятности задействованных рисков и  применить функцию предпочтения к этим распределениям вероятности.   Полученная мера риска должна удовлетворять желаемым свойствам. Эти свойства включают в себя неприятие риска (что является основополагающим для страхования) и предпочтение разнообразия (что  является основополагающим для теории портфелей и выбора инвестиций). \\
Для того, чтобы назначить цену или премию , нужно перевести случайную будущую прибыль или убыток  в денежный эквивалент. Для этого ,  кроме распределения вероятности этой прибыли или убытка, требуется также принцип назначения премии или цены.  Таким образом, компонентами модели для оценки рисков являются статистическая модель для распределения рисков, экономическая модель для предпочтения рисков и принцип назначения премий для преобразования меры риска в денежный эквивалент.\\
Цены или премии  также должны удовлетворять некоторым базовым свойствам. Используемая модель должна предоставлять рациональные и  непротиворечивые   результаты. Это включает в себя соответствие с наблюдаемым поведением  на финансовом или страховом рынке и рациональное обращение с разными видами рисков, страховыми и инвестиционными. \\
В актуарной науке было предложено множество принципов назначения премий (например, принцип отсутствия арбитража или принцип пропорционального деления функции риска).   Отметим, что принцип назначения премий согласуется с поиском детерминированного эквивалента риска в двойственной теории предполагаемой полезности, разработанной Яри. Эти подходы к оценке страховых контрактов обращаются со  страховыми потерями как с положительными случайными величинами и назначают премии выше, чем предполагаемое значение страховых потерь.

\section {Основные результаты}

\subsection{  Допущения и обозначения}
Случайные прибыль и потери, возникающие из страховых потерь и инвестиционных решений, обозначаются случайными величинами $X_i.$ Для каждой случайной прибыли или потери обозначим функцию распределения как $F_i(x)=P(X_i \leq x) $ и дополнительную функцию распределения как ${\overline F_i}(x)=P(X_i > x).$ Обобщенная обратная к дополнительной функции распределения определяется как 
$${\overline F^{-1}}(q) = inf\{x: {\overline F}(x) \leq q\}, 0 \leq q < 1, {\overline F^{-1}}(1)=0$$

Мы смотрим на все с точки зрения человека-страхователя, поэтому убыток $X_i$ является неположительным.  Выплата, совершаемая страховщиком, когда происходит страховой случай, обозначается $Y_i$ и является неотрицательной. Обычно страховая компания будет покрывать  убытки от изменения стоимости активов, происходящие  (только) из заранее оговоренных событий. В этом случае человек-страхователь одновременно является  держателем  этих активов и будет подвержен инвестиционным рискам, проистекающим из колебаний, вызванных экономическими и другими факторами. Полный риск по активу включает как экономические колебания в цене, так и потери по застрахованным событиям.\\
Мы предполагаем, что любая модели оценки рисков и назначения цен и премий должна удовлетворять свойству, что функция предпочтения рисков должна удовлетворять "неприятию риска" и "разнообразию портфеля".  Оба эти свойства считаются обязательными для финансовой модели, используемой для оценки рисков и распределения активов. Многие реальные финансовые решения, например, покупка страхования и инвестиционные решения, соответствуют  свойствам  "неприятия риска" и "разнообразия портфеля". Разнообразие портфеля подразумевает неприятие риска для предполагаемых моделей применения. Однако свойства неприятия риска недостаточно для наличия свойства разнообразия портфеля в случае общих функций предпочтения рисков (например, если она не будет выпукла вниз). \\
Мы определяем неприятие риска и разнообразие портфеля следующим образом:\\

{\bf Определение 1} { \itshape Неприятие риска }выражается в том, что неслучайная сумма (например, матожидание потери по риску) предпочтительнее, чем риск.\\
{\bf Определение 2} { \itshape Предпочтение разнообразия портфеля } выражается в том, что выпуклая комбинация (то есть линейная комбинация точек, где все коэффициенты неотрицательны, и их сумма равна 1) предпочтительнее, чем один риск, в предположении, что риски эквивалентны.\\

Предполагается, что на функциях распределения существует отношение предпочтительности $ \succeq$,  где символ $\succ$ обозначает строгую предпочтительность, а символ $\sim$ обозначает безразличие/эквивалентность. Например, выражение $F_1 \succ F_2$ обозначает, что случайные прибыль или потери  $X_1$ с функцией распределения $F_1$  строго предпочтительнее, чем случайные прибыль или потери  $X_2$ с функцией распределения $F_2$.\\

\subsection {Ожидаемая полезность и нейтральная мера}
Аксиоматический подход к выводу функции предпочтения, предназначенной для упорядочивания рисков, использующий ожидаемую полезность, был дан фон Нейманом и Моргенштейном. Аксиомы также присутствуют у Ванга и Янга. Ключевая аксиомой является так называемая аксиома независимости. Эта аксиома утверждает, что если $X \succ Y$ и Z - произвольный риск, то
$$\{(\alpha,X),(1-\alpha,Z)\} \succ \{(\alpha,Y),(1-\alpha,Z)\} $$
для любых $0 \leq \alpha \leq 1 $, где $ \{ (\alpha,X),(1-\alpha,Z) \}$ - это вероятностная смесь с такой функцией распределения:
$$F_{  \{(\alpha,X),(1-\alpha,Z)\} } = \alpha F_X(x) + (1-\alpha)F_Z(x),$$
что эквивалентно:
$$\overline F_{  \{(\alpha,X),(1-\alpha,Z)\} } = \alpha {\overline  F}_X(x) + (1-\alpha ){ \overline F} _Z(x)$$.

Мы это можем переписать как 
$$\alpha {\overline  F}_X(x) + (1-\alpha ){ \overline F} _Z(x) \succ \alpha {\overline  F}_Y(x) + (1-\alpha ){ \overline F} _Z(x).$$

Предположим, что у человека есть изначальный капитал W. Из этой аксиомы можно показать, что существует функция полезности u такая что 
$F_1 \succ F_2 \Leftrightarrow  E \left[ u(W + X_1) \right]  >  E \left[u(W + X_2) \right] .$  

Действительно, докажем существование функции полезности, согласованной с данным отношением предпочтения:\\


{\large \bfТеорема.} Пусть отношение предпочтения полно (то есть любые два блага сравнимы) и транзитивно, причем множество благ конечно. \\
Тогда существует такая функция предпочтения $u(x): X \rightarrow R$, что\\
 $u(x) \geq u(y) \leftrightarrow x \succeq y$

{\large \bf Доказательство.} Для фиксированного $x$, обозначим \\
$M(x) := \{ y \in X | x \succeq y\}$\\
И положим $u(x):= |M(x)|$\\
Проверим, что это и есть искомая функция полезности:\\
1) Пусть $x \succeq y$. Возьмем произвольное $z \in M(y)$, тогда по определению $M(y)$ имеем: $y \succeq z$.\\
В силу полноты, $z$ и $x$ сравнимы. Из транзитивности получаем, что $x \succeq z$, то есть $z \in M(x)$.
То есть мы показали, что $M(y)$ вложено в $M(x)$, то есть мощность у $M(x)$ больше. 
То есть $u(x) = |M(x)| \geq M(y) = u(y)$\\
2) Пусть $u(x) \geq u(y)$. Из полноты мы знаем, что либо $x \succeq y$, либо $y \succeq x$. 
Тогда по первому пункту получим либо $M(y)$ вложено в $M(x)$, либо $M(x)$ вложено в $M(y)$ соответственно. То есть $M(x)$ и $M(y)$ не могут не пересекаться или пересекаться только частично. Но по условию мощность $M(x)$ больше, чем мощность $M(y)$, поскольку $u(x) \geq u(y)$.
 В силу рефлексивности  $y \in M(y)$, поэтому  $y \in M(x)$, то есть $ x \succeq y$. Чтд. \\


Поскольку мы предполагаем неприятие риска, $u(x)$ - это возрастающая выпуклая вверх функция от X,  которая дифференцируема по крайней мере 2 раза, причем $u^{'} > 0$ и $u^{''} < 0.$ Функции полезности единственны с точностью до аффинных преобразований, то есть $u^{*}(X) = au(X) + b $ дает то же самое упорядочивание рисков, что и $u(X).$ Функции полезности можно стандартизовать, положив $u(k)=0$ и $u^{'}(k) = 1$ для некоторой точки k.\\
Для страхования убыток $X_i$ будет отрицательным. Если мы запишем потери как положительную случайную величину $Y=-X,$ то $ E \left[ u(X) \right]  = E  \left[ u(-Y) \right] . $\\

Имеем $$ E \left[  u(X) \right]   \leq  u(E  \left[  X \right]  ),$$
 из-за неприятия риска ( дейтвительно, $u(E  \left[  X \right]  ) >  u(X)$, а потом взять матожидание от обеих частей,  но на правой части это никак не отразится, поскольку там стоит число), и это верно для всех рисков, где X отрицательно для страховых потерь.\\
Аналогично (просто взяв монотонно возрастающую функцию u от обеих частей предыдущего неравенства) мы имеем $$u^{-1}  \left[ E \left[ u(X) \right] \right]   \leq  E  \left]  X \right]  $$

Для страховых рисков, обозначив $Y=-X,$ где Y - это неотрицательная случайная величина, получим 
$$ - u^{-1}  \left[ E \left[ u(X) \right] \right]   \geq  E  \left]  X \right]  $$

Безрисковый эквивалент  определяется так:\\

{\bf Определение 3}  Если человеку безразлично обладание риском и получение неслучайной суммы, то эта неслучайная сумма называется { \itshape безрисковым эквивалентом } этого риска.\\

Обозначим безрисковый эквивалент риска X как $C_X.$ Для человека с текущим капиталом W мы имеем:
$$u(W + C_X) = E \left[   u(W + X)  \right]  $$
$$\Leftrightarrow C_X = u^{-1}  \left[  E \left[   u(W + X)  \right]  \right]  - W$$

Для страховых рисков, обозначив величину убытка $Y=-X$ и безрисковый эквивалент как $\pi_Y,$ то

$$ u(W  -  \pi_Y ) = E \left[   u(W - Y)  \right] $$
$$\Leftrightarrow  \pi_Y =  W- u^{-1}  \left[  E \left[   u(W - Y)  \right]  \right]  $$

Подчеркнем, что $  \pi_Y = - C_{-Y}$

По неравенству Йенсена для любой случайной величины X и выпуклой вверх функции u:
$$ u(W + E \left[ X \right] )  \geq  E \left[   u(W + X)  \right]  $$

Поэтому из системы 
$$\left[ \begin{array}{crl}
u(W + C_X) = E \left[   u(W + X)  \right] \\
u(W + E \left[ X \right] )  \geq  E \left[   u(W + X)  \right]\\
X \succeq Y \Leftrightarrow u(X) \geq  u(Y)
\end{array}\right. $$

получаем:

$$\Rightarrow  E \left[ X \right]   \geq  C_X.$$

Аналогично получаем  $$  \pi_Y \geq  E \left[ Y \right] $$ для неотрицательной случайной величины потерь Y.\\

Рассмотрим портфель рисков $ \{ X_1, \dots, X_n \} ,$ заданный как $\sum\limits_{i=1}^{n} \alpha_i X_i.$ Из-за свойства предпочтения разнообразия портфеля мы считаем, что человек предпочитает владение таким портфелем рисков владению любым из рисков портфеля по отдельность, в предположении идентичности всех этих рисков $X_i.$ Поэтому из-за предпочтения разнообразия портфеля мы имеем:
$$C_{portfolio} \succ C_{X_i}$$
где $C_{X_1}= C_{X_2} = \dots = C_{X_n}$\\

Для ожидаемой полезности свойство предпочтения разнообразия тоже выполнено.\\

Для ценообразования, рассмотрим такой выбор портфеля рисков $\{ X_i: i=1, \dots, n \}$ и страховых полисов с положительными выплатами  по убыткам $\{ Y_i: i=n+1, \dots, m \}, $ чтобы максимизировать ожидаемую полезность богатства. Предположим, что риск (или полис) номер i имеет цену $P_i.$ Для инвестиционных рисков $X_i$: $P_i$ является величиной начальных вложений, а для страховых полисов с выплатами $Y_i$: $P_i$ является премией.\\

Наша задача максимизировать 
$$E \left[ u \left(  \sum\limits_{i=1}^{n} \alpha_i X_i +  \sum\limits_{i=n+1}^{m} \alpha_i Y_i \right)  \right] $$
с учетом того, что $$W=  \sum\limits_{i=1}^{m} \alpha_i P_i$$

То есть у нас изначально было W денег, мы на часть денег купили рискованные финансовые инструменты, а на оставшуюся часть денег приобрели договор страхования.

Дальше нам понадобится метод множителей Лагранжа с ограничениями типа равенство, поэтому напомним его. \\
Рассмотрим следующую задачу на условный экстремум:
$$\left[ \begin{array}{crl}
\phi_{0} (x) \rightarrow extr\\
\phi_{i} (x)=0 , i=1, ..., m.\\
\end{array}\right. $$

Допустимое множество $\Omega$ имеет вид
$$\Omega = \{x \in R^{n} | \phi_{i}(x)=0, i =1, ..., m\}$$

Будем считать, что все функции $\phi_{i}, i = 0,1 ... m$ непрерывно дифференцируемы. Тогда задача называется гладкой конечномерной задачей с ограничением типа равенство. Для решения таких задач определим функцию Лагранжа  от пары переменных $x,\lambda$\\
$L(x,\lambda) = \sum\limits_{i=0}^{m} \lambda_{i} \phi_{i}(x), \lambda = (\lambda_{0}, \lambda_{1}, ... \lambda_{m}) \in R^{m+1}$\\
При этом координаты вектора $\lambda$ называют множителями Лагранжа.\\
Сформулируем необходимое условие экстремума в этой задаче.\\

{\large \bf Теорема}(принцип Лагранжа) Если $a \in \Omega $ является точкой условного экстремума данной задачи, то существует такой ненулевой вектор $\lambda \in R^{m+1}$, что $L^{'}_{x} (a, \lambda) = 0$. То есть необходимым условием существования экстремума является разрешимость системы из $n+m$ уравнений, полученных дифференцированием функции Лагранжа по $x_i$ и $\lambda_j$.\\


Применим метод множителей Лагранжа для нахождения условного экстремума нашей функции (условие - наше единственное ограничение  $W=  \sum\limits_{i=1}^{m} \alpha_i P_i$)

Функция Лагранжа  равна
$$  L(\alpha_1, \dots, \alpha_m, \lambda)= E \left[ u \left(  \sum\limits_{i=1}^{n} \alpha_i X_i +  \sum\limits_{i=n+1}^{m} \alpha_i Y_i \right)  \right]  - \lambda  \left[ W - \sum\limits_{i=1}^{m} \alpha_i P_i   \right] $$



Дифференцируем $L(\alpha_1, \dots, \alpha_m, \lambda)$ по $\alpha_i$ и приравниваем производные к нулю:

$$ L^{'}_{\alpha_i}= E \left[  u^{'}(W^{*}) X_i \right]  + \lambda P_i = 0  ,   i = 1, \dots, n$$

$$ L^{'}_{\alpha_i}= E \left[  u^{'}(W^{*}) Y_i \right]  + \lambda P_i= 0  ,  i = n+1, \dots, m$$

где $ W^{*} = \sum\limits_{i=1}^{n} \alpha^{*}_i X_i +  \sum\limits_{i=n+1}^{m} \alpha^{*}_i Y_i $\\

Следовательно, 
$$  \lambda P_i = -  E \left[  u^{'}(W^{*}) X_i \right] ,   i = 1, \dots, n$$
$$  \lambda P_i = -  E \left[  u^{'}(W^{*}) Y_i \right] ,   i = n+1, \dots, m$$


Предположим, что $X_i$ не содержит риска обладания инвестиционным или страховым риском. То есть это безрисковое владение текущим капиталом без подверженности экономическим рискам или страховым событиям. Тогда $$ P_i= X_i= k,$$  где k - константа, тогда оптимальное значение лагранжева множителя $\lambda$ есть 

$$ \hat  {\lambda } = -  E \left[  u^{'}(W^{*}) \right] $$

Тогда мы имеем 

$$  P_i = \frac { E \left[  u^{'}(W^{*}) X_i \right] } { E \left[  u^{'}(W^{*}) \right] },   i = 1, \dots, n$$
$$  P_i = \frac { E \left[  u^{'}(W^{*}) Y_i \right] } { E \left[  u^{'}(W^{*}) \right] },   i = n+1, \dots, m$$


Следовательно, мы можем выразить цену/премию за риск $X_i$ как 
$$P_i= E \left[ \Psi X_i  \right]$$

а премию за страховой контракт как 

$$P_i= E \left[ \Psi Y_i  \right]$$

где $$ \Psi =  \frac 
{  u^{'}(\sum\limits_{i=1}^{n} \alpha^{*}_i X_i +  \sum\limits_{i=n+1}^{m} \alpha^{*}_i Y_i )} 
{ E \left[  u^{'}(\sum\limits_{i=1}^{n} \alpha^{*}_i X_i +  \sum\limits_{i=n+1}^{m} \alpha^{*}_i Y_i ) \right] }$$


Теперь используем случайную величину $\Psi$ как производную Родона-Никодима для того, чтобы изменить вероятностную меру P на Q так, чтобы для новой меря мы имели $P_i= E_{Q} \left[ X_i  \right].$
Тогда для новой вероятностной меры ценой каждого риска является его матожидание, поэтому мы можем называть меру Q  "нейтральной мерой риска". \\
Премии за страховые контракты, использующие такое преобразование меры, являются аддитивными для всех рисков. Поэтому они соответствуют нашему требованию непротиворечивого обращения с рисками (ну было бы логично, что премии  можно складывать, и так и получилось).\\


\subsection { Двойственная теория полезности и функция искажения }

Двойственная теория Яри разработала альтернативу ожидаемой полезности, используя аксиоматический подход, в котором аксиома независимости заменена на двойственную аксиому независимости.  Двойственная аксиома независимости говорит, что если $ X \succ Y$ и Z- произвольный риск, то

$$\alpha {\overline  F}^{-1}_X(x) + (1-\alpha ){ \overline F}^{-1} _Z(x) \succ \alpha {\overline  F}^{-1}_Y(x) + (1-\alpha ){ \overline F}^{-1} _Z(x).$$
для любых $0 \leq \alpha \leq 1.$

Эти аксиомы подразумевают, что существует такая непрерывная неубывающая функция g, что 

$$F_1 \succ F_2 \Leftrightarrow 
 - \int^{1}_{0} g(q)d{\overline  F}^{-1}_1(q)   >   - \int^{1}_{0} g(q)d{\overline  F}^{-1}_2(q)  $$ 


Подставив $q={\overline  F}_i(x),$ получим,  что:

$$  - \int^{1}_{0} g(q)d{\overline  F}^{-1}_1(q)  = \int^{\infty}_{0} g({\overline  F}_i(x))dx$$

Яри отмечает, что $U_{X_i} =   \int^{\infty}_{0} g({\overline  F}_i(x))dx$ является полезностью, которая приписывает случайной величине безрисковый эквивалент. Таким образом, человек будет безразличен по отношению к получении неслучайной суммы $U_{X_i}$ или риска $X_i.$\\

Ванг предлагает назначать цену страховым рискам, используя функцию искажения, основанную на пропорциональном преобразовании риска. Для страхового риска $Y$ (неотрицательной случайной величины) с дополнительной функцией распределения $ {\overline  F}_Y(x) = P(Y > x),$ Ванг предлагает принцип назначения премий $H_{r}(X) = \int^{\infty}_{0} ({\overline  F}_i(x))^{r}dx, $ где $0 \leq r \leq 1.$

 $H_{r}(X)$ используется для расчета премий с поправкой на риск. Отметим, что функция $g(x)= x^{r}, 0 \leq r \leq 1$ выпукла вверх.\\
 
 Ванг и Янг определяют функцию искажения как неубывающую функцию, удовлетворяющую условиям  $g(0)=0$ и $g(1)=1,$  такую что для неотрицательной случайной величины Y безрисковые эквиваленты  
 $ H_{g} \left[ Y \right] $ и $ H_{g} \left[ -Y \right] $   задаются как 
 
 $ H_{g}(Y) = \int^{\infty}_{0} g({\overline  F}_Y(x))dx 
 =  - \int_{0}^{1} g(q)d{\overline  F}^{-1}_Y(q) \\
 = - (g(q){\overline  F}^{-1}_Y(q) \|_{0}^{1}) + \int_{0}^{1} {\overline  F}^{-1}_Y(q))dg(q)
 =  \int_{0}^{1} {\overline  F}^{-1}_Y(q))dg(q)$\\
 
 
 и $$ H_{g} (-Y) = - H_{\tilde g} (Y),$$  
 
 где $\tilde g $ - это функция искажения, определяемая как $ \tilde g = 1 - g(1-q), 0 \leq r \leq 1.$\\
 
 Отметим, что  $\tilde g $ выпукла вниз, если g выпукла вверх.\\
 
   \begin{figure}[h!]
\center{\includegraphics[scale=0.9]{pictNEW.pdf}}
\caption{   Графики $g(t), \tilde g (t)  $ }
\label{fig:image}
\end{figure}

 
 
 
 Безрисковый эквивалент  $ H_{g}(Y) $ удовлетворяет условиям:\\
 
$\bullet$ Если g выпукла вверх, то  $ H_{g}(Y)  \geq E(Y).$\\
$\bullet$  $ H_{g}(aY+b) = a H_{g}(Y) + b, a,b \geq 0.$\\
$\bullet$ Если g выпукла вверх, то $ H_{g}(Y_1 + Y_2 )  \leq  H_{g}(Y_1)  + H_{g}(Y_1).$\\
 

Третье свойство называется суб-аддитивность. $H_{g}(Y)$ будет аддитивна в специальном случае 
ко-монотонных рисков. Риски $X_1$ и  $X_2$ являются ко-монотонными, если существует риск Z и неубывающие действительнозначные функции f и h, такие что $X_1= f(Z)$ и  $X_2= h(Z).$  Понятие ко-монотонных рисков является расширением полной корреляции. \\

Если бы нам нужно было применить этот безрисковый эквивалент к инвестиционным рискам, то из-за неприятия риска, мы бы требовали, чтобы $H_{g} \left[ X \right] \leq E \left[ X \right] , $ что означало бы, что g выпукла вверх. Для страхового риска $Y=-X$ (неотрицательной случайной величины), согласно Вангу и Янгу мы имеем:

$$H_{g} \left[ Y \right]  =  H_{g} \left[ - X \right]  = - H_{\tilde g}(X) \eqno (1)$$

В этом случае $\tilde g$ выпукла вниз из-за неприятия риска, поэтому g тоже выпукла вверх. Поэтому дальше мы получаем, что 
$$H_{\tilde g}(X) \leq E \left[ X \right]$$

$$\Rightarrow - H_{\tilde g}(X) \geq  - E \left[ X \right] =  E \left[ Y \right]$$

$$\Rightarrow H_{ Y}(X) \geq   E \left[ Y \right]$$

Отметим, что для выпуклой вверх функции g:

$$H_g \left[ \sum\limits_{i=1}^{n} \alpha_i X_i   \right]\geq  H_{g} \left[ Y \right] ,$$

если $ H_g \left[ X_i \right] = H_g \left[ X_1 \right] , i= 1 \dots, n.$ Таким образом, свойство разнообразия выполнено для инвестиционных рисков.\\

Если мы рассматриваем страховые убытки как неотрицательные случайные величины $Y_i$, то необходимо использовать выпуклую вверх функцию g, соответствующую выпуклой вверх функции ${\tilde g},$ используемой для инвестиционного риска $X_i.$ Однако свойство разнообразия не выполняется для выпуклой функции g. В результате,  рисковая мера искажений  не подходит для управления активами и пассивами. Пассивы (или убытки по рискам) и риски  активов обрабатываются не  соответствующим образом, поскольку для рисков  активов  выполнено свойство разнообразия, но оно не выполнено для рисков пассивов/обязательств. \\

 Мы уже отмечали, что безрисковый эквивалент  $ H_{g}(Y) $ аддитивен только для ко-монотонных рисков. Но мы требуем, чтобы принцип назначения премий был аддитивным, поскольку цены финансовых активов аддитивны.  В действительности означает, что страховые премии назначаются на рынке  на основе полной информации, без учетов стоимости транзакций и других дефектов, а также с учетом принципа отсутствия арбитража на страховом рынке. Это соответствует понятию равновесия на страховом ранке и ценообразованию на финансовые активы при идеальных рыночных условиях. В этой статье мы не будем разбирать задачи ценообразования при ослаблении этих условий.\\
 
 Для меры риска   $ H_{g} \left[ X \right]  $ при условии выпуклости вверх функции g мы имеем свойства неприятия риска и разнообразия портфеля. Но мы не хотим использовать этот принцип, поскольку мы хотим, чтобы премии были аддитивными. Это свойство аддитивности выполняется только для ко-монотонных рисков, а мы хотим, чтобы оно выполнялось для всех рисков. \\
 
Заинтересованные подходом функцией искажения  в построении  меры риска $ H_{g} \left[ X \right]  ,$ мы предлагаем метод назначения премий, кторый будет обладать свойством аддитивности премий. Кроме того, на основе этой меры ценообразования мы определяем безрисковый эквивалент, используя выпуклые вверх функции, похожие на функции полезности. Этот безрисковый эквивалент обладает желаемыми свойствами неприятия риска, разнообразия портфеля и еще имеет более логичное упорядочивание рисков, похожее на упорядочивание функцией полезности.


\subsection { Изменение меры для премий с использованием функций искажения }

Подход  Ванга к страховым премиям  с использованием функций искажения  использует распределение потерь как положительные случайные величины. Получающаяся премия эквивалента безрисковому эквиваленту в дуальной теории ожидаемой полезности.  Если бы нам нужно было применить этот безрисковый эквивалент к инвестиционным рискам, то мы бы использовали выпуклую вверх функцию $\tilde g,$  эквивалентную выпуклой вверх функции g, которая используется для страхового ценообразования.\\

Страховая премия $\geq$ матожидания риска, это соответствует свойству неприятия риска. Однако, если наш подход  функций искажения  для страховых рисков применить для их ценообразования, то получающиеся страховые премии оказываются неаддитивными, кроме случая ко-монотонных рисков.\\

Мы требуем такой принцип назначения премий для портфеля рисков, что он удовлетворяет свойству неприятия риска, и соответствующая мера риска не противоречит использованию выпуклой вверх функции полезности, которая применяется к страховым рискам.\\

Для выплат по портфелю $Y_\alpha = \sum\limits_{i=1}^{n} \alpha_i Y_i $ страховых контрактов, где выплачиваемые суммы являются положительными случайными величинами $\{ Y_1, \dots, Y_n \}$ с распределениями $(F_1, \dots, F_n),$ мы назначаем премию 
$$\pi_{r}(Y_\alpha) = E_{P_{n, r}}Y_{\alpha},$$ 
где $P_{n,r} $ есть вероятностная мера, соответствующая распределению
$$ F_{n,r} = \prod_{i=1}^{n} F_{i}^{r}, 0 < r \leq 1.$$

Очевидно, что такой принцип назначения премий будет аддитивным по портфелю рисков, причем
$$\pi_{r}(Y_{\alpha}) \geq EY_{\alpha}.$$

Чтобы непротиворечиво обращаться с инвестиционными и страховыми рисками, мы определяем меру риска портфеля как 
$$ U_{r}(Y_{\alpha} = u^{-1}(E_{P_{n, r}}u(Y_{\alpha})),$$
где u выпукла вверх.\\

{\bf  Теорема 1} Пусть u - выпуклая вверх, возрастающая и дважды дифференцируемая нелинейная функция, $u^{'} > 0,$  и пусть матожидания $E_{P_{n,r}}Y_{j}$ - непрерывные слева функции от r в точке $r=1.$\\
Тогда существует такое $0 < r^{*} < 1,$ что для всех $r^{*} \leq r < 1$ и любого $\alpha = (\alpha_1, \dots, \alpha_n),$  причем $\sum\limits_{i=1}^{n} \alpha_i = 1,$ будет верно:

$$U_{r}(Y_{\alpha} ) \leq EY_{\alpha},$$

то есть мера риска $U_{r}( \dot ) $ обладает свойством неприятия риска.\\

{\bf  Доказательство} Из выпуклости вверх  функции u следует, что $u(x) \leq u(0) + u^{'}(0)x.$ Без потери общности, gредположим, что $u(0) = 0.$ Тогда, из условий теоремы, функция 

$$ \Pi(\alpha_1, \dots, \alpha_n, r) = u^{-1} (  E_{P_{n,r}} u (\sum\limits_{i=1}^{n} \alpha_i Y_i))$$

является непрерывной слева функцией от r  в точке $r=1$ для любого $\alpha = (\alpha_1, \dots, \alpha_n).$ \\

Это значит, что $$\lim \limits_{r  \to 1-0}  \Pi(\alpha_1, \dots, \alpha_n, r) = \Pi(\alpha_1, \dots, \alpha_n, 1)  = u^{-1}(Eu  (\sum\limits_{i=1}^{n} \alpha_i Y_i)),$$

и эта сходимость равномерная по $ \alpha \in H = \{  \sum\limits_{i=1}^{n} \alpha_i = 1\},$ потому что гиперплосткость H компактна. \\

Тогда для $\Delta > 0$ существует $ 0 < r^{*} < 1$ такое что для всех $r^{*} \leq r < 1$ 

$$ \left| \Pi(\alpha_1, \dots, \alpha_n, r)  - u^{-1}(Eu  (\sum\limits_{i=1}^{n} \alpha_i Y_i)) \right| < \frac {\Delta}{2} \eqno (2)$$

С другой стороны, для выпуклой вверх нелинейной функции u, функция :

$$ \Delta(\alpha) = \pi_{r_{\left| r=1 \right.}} (Y_{\alpha} - \Pi(\alpha_1, \dots, \alpha_n, 1) $$ 

непрерывна по $\alpha$ для всех $\alpha \in H$\\

А поскольку $\alpha \in H,$ а H - компактное множество, то

$$ \inf \limits_{\alpha \in H} \Delta(\alpha) = \Delta(\alpha^{*}) > 0, \alpha^{*} \in H.$$

Тогда 

$$ EY_{\alpha} - U_{r}(Y_{\alpha})  =   \pi_{r_{\left| r=1 \right.}} (Y_{\alpha} - \Phi(\alpha_1, \dots, \alpha_n, \alpha )$$
$$ \geq \Delta(\alpha) - \left| \Phi(\alpha_1, \dots, \alpha_n, 1) - \Phi(\alpha_1, \dots, \alpha_n, \alpha ) \right|$$
$$ \geq \Delta(\alpha^{*}) - \Delta(\alpha^{*})/2 > 0, r^{*} < r < 1$$


Предложенная мера риска также будет обладать свойством разнообразия для страховых потерь, когда эти потри рассматриваются как отрицательные (неположительные) случайные величины.

\subsection{ Непротиворечивость меры риска для инвестирования и меры риска для потерь в отношениях предпочтения} 

Предположим, что мы непротиворечиво определяем меру риска $R_{m_{I}}(X) $ для инвестиций
и меру риска $R_{m_{L}}(X) $  для неотрицательных  рисков. В предыдущих разделах мы рассматривали некоторые специальные определения  $R_{m_{I}}(X) $ и $R_{m_{L}}(X) :$\\

1) Ожидаемая полезность: для выпуклой вверх функции u

$$ R_{m_{I}}(X)= u^{-1}  \left[  E \left[   u(W + X)  \right]  \right]  - W$$
$$ R_{m_{L}}(Y)= W- u^{-1}  \left[  E \left[   u(W - Y)  \right]  \right]  - W \eqno (3)$$

2) Функции искажения : для выпуклой вверх функции u

$$ R_{m_{I}}(X)= H_{\tilde g}(X)$$
$$ R_{m_{L}}(Y)= H_{g}(Y)$$

где $\tilde g (q) = 1 - g(1-q)$ - выпуклая вниз функция.

3) Изменение меры: для $r^{*} \leq r < 1$

$$  R_{m_{I}}(X) = U_r(X) = u^{-1}(E_{P_{n,r}}u(X))$$
$$  R_{m_{L}}(Y) =E_{P_{n,r}}(Y) \eqno (4)$$

Меры  $R_{m_{I}}(X) $ и $R_{m_{L}}(X) $ будут непротиворечивыми относительно отношения предпочтения, если они порождают одинаковое отношение предпочтения, то есть для двух неотрицательных рисков из портфеля:

$$ X \prec _{ R_{m_{I}}} Y \eqno (5) $$
 и 
 
 $$ X \prec _{ R_{m_{L}}} Y \eqno (6) $$

должны выполняться одновременно. Отметим, что меры, заданная как ожидаемая полезность, очень близки к нашему требованию непротиворечивости относительно отношения предпочтения. На самом деле, если мы определяем $R_{m_{I}}(X) $ и $R_{m_{L}}(X) $ по формулам (3), то из (5) мы получаем, что

$$ W + X \prec_{u} W + Y \eqno(7)$$,

где отношение предпочтения $\succ_{u}$ задается $Eu(\dot)$ и эквивалентно 

$$ W - Y \succ_{u} W - X \eqno(8)$$

Видно, что (7) и (8) могут считаться эквивалентными для любой функции u, симметричной относительно W, то есть $u(W+X) = -u(W-X).$\\

Мы также можем рассматривать менее ограничительные формы непротиворечивости мер для инвестиций и для потерь. Например, мы можем потребовать эквивалентность (5) и  (6) для инвестиций и потерь, которые упорядочены  во втором стохастическом доминировании (или порядке остановки потерь).

{\bf  Определение 4 } Риск X меньше, чем риск Y в SSD, если для всех $x \geq 0$

$$ \int_{x}^{\infty} \overline F_{X}(u)du \leq  \int_{x}^{\infty} \overline F_{Y}(u)du \eqno(9)$$

где $ \overline F_{X}(x) = 1 - F_{X}(x)$ - дополнительная функция распределения X.

Определение 4 эквивалентно тому, что $Eu(X) \leq Eu(Y).$



{\bf  Определение 5 }  Две меры $R_{m_{I}}(X) $ и $R_{m_{L}}(X) $ дают непротиворечивое отношение предпочтения по отношению к SSD, если для $X \prec_{SSD} Y$ соотношения (5)  и (6) выполнены одновременно. 

Покажем сначала, что вообще говоря, меры, построенные по функциям искажения, не дают  непротиворечивое отношение предпочтения по отношению к SSD. 


\subsection {Пример  }


  Рассмотрим $g(t) = 1 - (1-t)^m,$ $m>1,$ m- целое, тогда $\tilde g (t) = t^m.$
  
  
   \begin{figure}[h!]
\center{\includegraphics[scale=0.9]{pictNEW.pdf}}
\caption{   Графики $g(t), \tilde g (t)  $ }
\label{fig:image}
\end{figure}

  \begin{figure}[h!]
\center{\includegraphics[scale=0.9]{python_code1.pdf}}
\caption{   Графики $g(t), \tilde g (t)  $ }
\label{fig:image}
\end{figure}



Видно, что $g(t)$ выпукла вверх, а $\tilde g $ выпукла вниз, причем 

$$ R_{m_{I}}(X) = H_{\tilde g}(X) = \int_{0}^{\infty}  (\overline F_{X}(x))^m dx = E(min(X_1, \dots, X_m)) \eqno(10)$$

$$ R_{m_{L}}(Y) = H_{ g}(Y) = \int_{0}^{\infty}(1-  (F_{Y}(x))^m) dx = E(max(X_1, \dots, X_m)) \eqno(11)$$

где $(X_1, \dots, X_m)$ - независимые одинаково распределенные величины с функцией распределения $F_X(x),$ а  $(Y_1, \dots, Y_m)$ - независимые одинаково распределенные величины с функцией распределения $F_Y(x).$\\


 \begin{figure}[h]
\center{\includegraphics[scale=0.9]{pictures2.mps}}
\caption{   Графики $  \overline F_{X}(x), \overline F_{Y}(x)$ }
\label{fig:image}
\end{figure}



Предположим, что $$ \overline F_{X}(x) = I_{(-\infty, 0)} + (1-x)I_{\left[ 0,1 \right]}$$

$$ \overline F_{Y}(x) = I_{(-\infty, 1/4)} + (2-4x)I_{\left[ 1/4,1/2 \right]}$$

то есть $$X \sim R\left[ 0,1 \right]$$
$$Y \sim R\left[ 1/4,1/2 \right]$$





В этом случае, функция $$J(x) = \int_{x}^{1} (\overline F_{Y} (u) - \overline F_{X} (u))du$$

Мы видим, что она положительна на всем промежутке $\left[ 0,1 \right], $ а это значит, что $X \prec_{SSD} Y.$

С другой стороны, используя формулу матожидания min и max равномерных величин::

$$  E(min(X_1, \dots, X_m)) = a + \frac{1}{m+1}(b-a)$$

$$  E(max(X_1, \dots, X_m)) = b- \frac{1}{m+1}(b-a)$$

$$ \Rightarrow  R_{m_{I}}(X)  = E(min(X_1, \dots, X_m)) =  \frac{1}{m+1}  <  R_{m_{I}}(Y)  = E(min(Y_1, \dots, Y_m)) =  \frac{1}{4} + \frac{1}{4} \frac{1}{m+1}$$

$$ R_{m_{L}}(X)  = E(max(X_1, \dots, X_m)) = 1-  \frac{1}{m+1}  >  R_{m_{L}}(Y)  = E(max(Y_1, \dots, Y_m)) =  \frac{1}{2} - \frac{1}{4} \frac{1}{m+1}$$

Это показывает, что  $R_{m_{I}}(X) $ и $R_{m_{L}}(X) $ не являются непротиворечивыми по отношению к SSD для меры функций искажения.\\

Теперь рассмотрим меры риска $R_{m_{I}}(X) $ и $R_{m_{L}}(X), $  определенные соотношением (4), построенные заменой меры, предложенной в этой статье. Мы покажем, что они могут давать непротиворечивое отношение предпочтения относительно SSD после замены меры.\\


{\bf  Определение 6 }  $X \prec_{r-SSD} Y $ в r-SSD, если для всех $x \geq 0$

$$ \int_{x}^{\infty} \overline {F_{X}(u)^r }du \leq  \int_{x}^{\infty} \overline{ F_{Y}(u)^r} du \eqno(12)$$

Подставив $x=0$ и использовав (4), мы получим 

$$ R_{m_{L}}(X)  \leq R_{m_{L}}(Y) \eqno(13)$$

то есть  

$$ X \prec_{ R_{m_{L}}} Y  \eqno(14) $$

С другой стороны, из  (12) следует, что 

$$ E_{P_{n,r}}u(-X) \geq  E_{P_{n,r}}u(-Y)$$

$$\Rightarrow  -X \succ_{ R_{m_{I}}} -Y  \eqno(15) $$

Если u(x) нечетная, то из (15)  следует, что

$$X \prec_{ R_{m_{I}}} Y  \eqno(16) $$

и тогда $R_{m_{I}}(X) $ и $R_{m_{L}}(X) $  являются непротиворечивыми в предпочтении относительно r-SSD. Если же функция u не является нечетной, то мы имеем только (14) и (15).\\



{\bf  Замечание 1  }  Стохастическое доминирование второго порядка также определяется в экономической и финансовой литературе. \\
Риск Y предпочтительнее риска X, если  для всех $x \geq 0:$

$$ \int_{0}^{x} F_{Y}(u)du \leq  \int_{0}^{x} F_{x}(u)du \eqno(17)$$
 
 Если это выполнено, то (17) эквивалентно тому, что $Eu(X) \leq Eu(Y)$ для любой возрастающей выпуклой вверх функции u, то есть это свойство неприятия риска.\\
 
 
 {\bf  Замечание 2  }
 Отметим, что  $$  \int_{x}^{\infty} \overline F_{X}(u)du \leq  \int_{x}^{\infty} \overline F_{Y}(u)du \eqno (9)$$
  и 
  
  $$  \int_{0}^{x} F_{Y}(u)du \leq  \int_{0}^{x} F_{x}(u)du \eqno(17)  $$

 не эквивалентны.\\
 
 Например, в примере 1 
 
 $$F_X(x) = xI_{\left[ 0,1 \right]} + I_{(0,\infty)}$$
 
 $$F_Y(x) = (4x-1)I_{\left[ 1/4,1/2 \right]} + I_{(1/2,\infty)}$$
 
 Тогда функция  $$J(x) = \int_{0}^{x} (F_{X} (u) -  F_{Y} (u))du \eqno(18) $$
 не является знакопостоянной на отрезке $\left[ 0,1 \right].$\\
 
 Отметим, что меры, построенные по функциям искажения, не дают  непротиворечивое отношение предпочтения относительно SSD и в смысле нового определения SSD(см . (17))
 

\subsection {  Выводы  }

Страховые премии могут получаться из мер риска,  для которых отношение предпочтения удовлетворяет свойствам неприятия риска и разнообразия портфеля. В предположениях отсутствия арбитража на рынке и идеальных рыночных условиях, равновесные премии должны быть аддитивными. Кроме того, мы должны применять меру риска к обеим сторонам бухгалтерской ведомости и при этом получать непротиворечивое упорядочивание рисков.\\
Мера риска, основанная на искаженных вероятностях, не удовлетворяет свойству аддитивности равновесных премий для страховых премий.  Премии, основанные на искаженных вероятностях, обычно только суб-аддитивны, а аддитивными являются только для ко-монотонных рисков. Таким образом, для портфеля из разных классов страховых рисков, премии для разных классов будут противоречить свойству равновесности.\\
Мы предложили принцип назначения страховых премий, используя замену меры, которая делает премии по портфелю аддитивными. Этот принцип назначения премий также обладает свойствами неприятия риска, то есть премии всегда превышают ожидаемое значение потерь.  Кроме того, мы предложили меру безрискового эквивалента для риска, основываясь на выпуклой вверх функции полезности и замене меры, которая упорядочивает инвестиционные риски непротиворечиво относительно страховых рисков.

\subsection {  Литература  }
:
\newline  [1] G.I.Falin, On the optimal pricing of a portfolio of heterogeneous insurance risks, 2008,MSU
\newline [2] Landsman, Sherris, Risk measures and Insurance premium principles, Insurance: Mathematics and Economics
\newline [3] Zaks Y., Frostig E. and Levikson B. (2006) Optimal pricing of a heterogeneous portfolio for a given risk level. ASTIN Bulletin 36(1), 161-185. 
\newline [4] Halmos, P.R. (1974) Finite-Dimensional Vector Spaces. Springer.

\subsection{ Код картинок в METAPOST}

\begin{figure}[h]
\center{\includegraphics[scale=0.8]{code.pdf}}
\caption{Про квантиль"}
\label{fig:image}
\end{figure}











\end{document}



\end{document}
